% ACL 2027 submission
% Does Language Matter? Evaluating LLM-Agent Performance Across Languages
\documentclass[11pt]{article}
\usepackage[review]{acl}   % use [review] for anonymous submission; drop it for camera-ready
\usepackage{times}
\usepackage{latexsym}
\usepackage[T1]{fontenc}
\usepackage[utf8]{inputenc}
\usepackage{microtype}
\usepackage{graphicx}
\usepackage{booktabs}
\usepackage{amsmath}
\usepackage{multirow}
\usepackage{enumitem}
\usepackage{xcolor}
\usepackage{url}

\newcommand{\dn}[1]{\textcolor{red!70!black}{#1}}
\newcommand{\up}[1]{\textcolor{green!45!black}{#1}}

\title{Does Language Matter? \\ Evaluating LLM-Agent Performance Across Languages}

\author{Anonymous ACL submission}

\begin{document}
\maketitle

\begin{abstract}
LLM agents are increasingly deployed to carry out multi-step, tool-using work, yet they are
evaluated almost exclusively in English. We ask a simple question with a non-obvious answer:
when the \emph{same} task is expressed in a different language, does the agent behave
differently, and if so, \emph{where} in its pipeline?
We build a paired bilingual (English/Bangla) version of \textsc{OfficeBench}, holding the
environment, the tool schema and the checkers fixed while varying only the language of the
system prompt and the task instruction, and we run two open-weight models (31B and 20B
parameters) over 300 subtasks in each condition---1{,}200 agent executions in total.
We find that the effect of language is \textbf{model-dependent in direction}: the 31B model
loses $9.7$ points in Bangla ($p{<}0.001$, exact McNemar) while the 20B model \emph{gains}
$5.6$ points ($p{=}0.043$). As a consequence, model ranking is language-dependent: a gap that
is overwhelming in English ($p{=}5{\times}10^{-9}$) becomes statistically indistinguishable in
Bangla ($p{=}0.52$).
A failure-mode analysis explains the divergence. The 31B model's failures are
\emph{content-bound} and degrade when the instruction language changes; the 20B model's
failures are \emph{action-format-bound}---$96\%$ of its English failures involve a malformed
action---and the language shift reduces the prose leakage that causes them.
Finally, we show that a substantial share of the measured Bangla penalty is
\emph{evaluator-induced}: $26\%$ of the 31B model's Bangla failures are answers written in
Bangla script and rejected by English keyword checkers. Crediting them would raise its Bangla
score above its English score. We argue that multilingual agent evaluation requires
language-aware verification, and we release code, traces and failure codes.
\end{abstract}

\section{Introduction}

Language-model agents now plan over many steps, call tools, and edit real artefacts---
spreadsheets, calendars, documents, email \citep{yao2023react,liu2023agentbench,zhou2024webarena}.
The benchmarks that measure them, however, are written almost entirely in English
\citep{jimenez2024swebench,xie2024osworld,mialon2024gaia}. This is a consequential gap. Most
of the world does not work in English, and an agent that edits a payroll spreadsheet is
supposed to be usable by the person who owns the spreadsheet, in that person's language.

The usual assumption is that moving a task into a lower-resource language degrades performance,
because multilingual capability lags English in knowledge and reasoning benchmarks
\citep{ahuja2023mega,shi2023mgsm,bandarkar2024belebele,huang2023notall}. For \emph{agents} this
assumption has not been tested carefully, and it hides a more interesting question. An agentic
episode is not one forward pass. It is a loop of reading state, reasoning, emitting a
structured action, and consuming the result. Language could in principle perturb any of these
stages, and a benchmark score collapses all of them into a single bit of success.

We therefore ask three questions.
\textbf{RQ1:} Does expressing an equivalent task in a different language change end-to-end
agent success, and in which direction?
\textbf{RQ2:} Where does the change act---on task understanding and content, or on the
mechanics of tool interaction?
\textbf{RQ3:} How much of an apparent cross-lingual gap is a property of the \emph{agent}, and
how much is a property of the \emph{evaluation harness}?

To answer these we construct a strictly paired bilingual protocol over \textsc{OfficeBench}
\citep{wang2024officebench}, a benchmark of 300 office subtasks requiring one, two or three
applications. We translate only the system prompt and the task instruction into Bangla---a
language with roughly 240M speakers that is nevertheless under-represented in LLM training
data \citep{joshi2020state}---while holding the Docker environment, the file state, the tool
schema and the automatic checkers byte-identical across conditions. We evaluate two
open-weight models of comparable but distinct scale and lineage, and we code every failure.

Our findings overturn the simple ``lower-resource is worse'' expectation:

\begin{itemize}[noitemsep,topsep=2pt]
\item \textbf{Direction is model-dependent.} \texttt{gemma4:31b} falls from $44.0\%$ to
$34.3\%$ in Bangla ($p{<}0.001$); \texttt{gpt-oss:20b} \emph{rises} from $26.7\%$ to $32.3\%$
($p{=}0.043$). Aggregate ``multilingual degradation'' is not a property of the language.
\item \textbf{Rankings do not transfer.} The English gap between the two models is enormous
($p{=}5{\times}10^{-9}$); in Bangla it disappears ($p{=}0.52$). A leaderboard built in English
does not predict the ordering in another language.
\item \textbf{Failure mode predicts direction.} The models fail for different reasons. When
failures are format-bound, a language shift can help; when they are content-bound, it hurts.
\item \textbf{Evaluators are part of the measurement.} $26\%$ of the 31B model's Bangla
failures are Bangla-script answers rejected by English keyword checks. The measured language
penalty is, in large part, a property of the checker.
\end{itemize}

\section{Related Work}

\paragraph{Agent benchmarks.} Interactive benchmarks evaluate LLMs as decision makers rather
than text producers: \textsc{AgentBench} \citep{liu2023agentbench}, \textsc{WebArena}
\citep{zhou2024webarena}, \textsc{OSWorld} \citep{xie2024osworld}, \textsc{SWE-bench}
\citep{jimenez2024swebench} and \textsc{GAIA} \citep{mialon2024gaia}. \textsc{OfficeBench}
\citep{wang2024officebench} is the setting we adopt: it requires switching between office
applications and grounding actions in a real file system, and it scores outcomes with
execution-based checks rather than string overlap. All of these are English-only.

\paragraph{Multilingual evaluation.} Cross-lingual evaluation of LLMs is well developed for
static tasks---\textsc{XNLI} \citep{conneau2018xnli}, \textsc{MEGA} \citep{ahuja2023mega},
\textsc{MGSM} \citep{shi2023mgsm}, \textsc{Belebele} \citep{bandarkar2024belebele}---and
consistently reports an English advantage, with the gap widening for lower-resource languages
\citep{huang2023notall,zhang2023dont}. \citet{etxaniz2024multilingual} show that models often
perform better when a prompt is first translated into English, implying an English-centric
internal pivot. These studies evaluate single-turn outputs; we ask what happens when the output
must be a \emph{valid action in a loop}.

\paragraph{Tool use and structured output.} Agents must emit well-formed, schema-conforming
calls \citep{schick2023toolformer,qin2024toolllm,patil2024gorilla}. Format adherence is itself a
capability, and it is known to be brittle. Our analysis shows that this brittleness, rather than
comprehension, dominates one of our two models and determines the sign of its language effect.

\paragraph{Failure analysis.} Rather than reporting a single score, a growing line of work
dissects \emph{how} systems fail \citep{bang2023multitask}. We follow that approach: we code
every failed episode and treat the code distribution, not the aggregate, as the object of study.

\section{Study Design}

\subsection{Benchmark and environment}
\textsc{OfficeBench} provides 300 subtasks spanning eight applications (calendar, excel, ocr,
pdf, shell, word, email, and an in-container \texttt{llm} helper). Tasks are stratified by the
number of applications required: 93 one-app, 95 two-app and 112 three-app subtasks. Each
subtask ships a seeded file system and a set of execution-based checks (cell values, file
existence, calendar overlap, keyword containment). A subtask passes only if \emph{all} of its
checks pass. Every episode runs in a fresh Docker container reset from a pinned image.

\subsection{The paired bilingual protocol}
Our central design requirement is that language be the \emph{only} variable. Concretely, across
the two conditions we hold fixed: the container image and seeded files; the available
applications; the tool schema (action names and argument keys); the step budget; the decoding
parameters; and the checkers. We vary:

\begin{enumerate}[noitemsep,topsep=2pt]
\item \textbf{The system prompt}, authored natively in Bangla, including the role description,
the application descriptions and the response-format instruction.
\item \textbf{The task instruction}, machine-translated into Bangla and cached so that every
model sees an identical string for a given subtask.
\end{enumerate}

Crucially, the \textbf{tool schema remains in English}. Action identifiers such as
\texttt{switch\_app} and argument keys such as \texttt{file\_path} are API symbols that the
executor matches literally; translating them would change the environment rather than the
language of the task, and would make the two conditions incomparable. We therefore present the
schema in English in both conditions and attach a Bangla gloss describing what each tool does.
This isolates \emph{instruction language} from \emph{interface language}, and it mirrors
deployment reality, where an interface is typically English while the user is not.

\subsection{Agent scaffold}
We use a ReAct-style loop \citep{yao2023react}. At each step the model receives the system
prompt, the task, the history of previous (action, observation) pairs, and the actions
available in the current application. It must reply with a short reasoning segment followed by
exactly one JSON action:

\begin{quote}\small\ttfamily
\#\#Thought: \textless reasoning\textgreater \\
\#\#Command: \{"app": ..., "action": ..., ...\}
\end{quote}

The harness parses the JSON after the last \texttt{\#\#Command:} marker, so reasoning may not
corrupt the action. An episode terminates when the agent emits \texttt{finish\_task}, repeats an
identical action five times (recorded as \emph{stuck}), or reaches the 20-step budget. We log the
full trajectory---prompt, reasoning, action, observation---for every step.

\subsection{Models}
We evaluate two open-weight models served through the same API with identical decoding
($\mathrm{temperature}=0$, $\mathrm{top\_p}=1$): \texttt{gemma4:31b} ($32.7$B parameters) and
\texttt{gpt-oss:20b} ($20.9$B parameters).\footnote{Served via Ollama; sizes read from the
served model metadata.} The two are deliberately different in kind: the former is a
conventional instruction-following model, the latter exposes a separate reasoning channel, and
this architectural difference turns out to drive our central result. For the reasoning model we
fold the reasoning channel into the logged \texttt{\#\#Thought} so that both conditions are
logged identically.

\subsection{Metrics and statistics}
Our primary metric is the benchmark's own pass rate. Because the design is \emph{paired}---the
same 300 subtasks in both conditions---we test differences with the exact two-sided McNemar
test over discordant pairs, which is the appropriate test for paired binary outcomes
\citep{mcnemar1947}. We report discordant counts ($b$, $c$) alongside each $p$-value so that
effect size is visible.

\subsection{Failure coding}
We code every failed episode from its trajectory. Besides a general taxonomy (benchmark defect,
wrong content, out of steps, premature finish, task read differently, action format), the Bangla
condition requires codes that cannot occur in English, notably \emph{answer language}: the agent
produced a semantically reasonable answer in Bangla script which an English keyword check
rejects. In addition to coding, we extract fully objective signals for all 1{,}200 episodes:
whether any action was rejected as malformed, whether the episode ended stuck, whether it
exhausted the step budget, whether it finished and still failed, episode length, and whether the
final answer contained Bangla script.

\section{Results}

\subsection{RQ1: Language changes success, but not in one direction}

Table~\ref{tab:main} gives the main result. The two models move in \emph{opposite} directions.
\texttt{gemma4:31b} drops from $44.0\%$ to $34.3\%$, a $9.7$-point loss, with 41 subtasks
passing only in English against 12 passing only in Bangla (exact McNemar $p=8.2\times10^{-5}$).
\texttt{gpt-oss:20b} moves the other way, from $26.7\%$ to $32.3\%$, a $5.6$-point gain, with 23
English-only passes against 40 Bangla-only passes ($p=0.043$).

\begin{table}[t]
\centering\small
\begin{tabular}{llcccc}
\toprule
Model & Lang & 1-app & 2-app & 3-app & All \\
 & & (93) & (95) & (112) & (300) \\
\midrule
\multirow{2}{*}{gemma4:31b}
 & EN & 66.7 & 54.7 & 16.1 & \textbf{44.0} \\
 & BN & 45.2 & 45.3 & 16.1 & \dn{34.3} \\
\midrule
\multirow{2}{*}{gpt-oss:20b}
 & EN & 44.1 & 34.7 & 5.4 & 26.7 \\
 & BN & 53.8 & 38.9 & 8.9 & \up{32.3} \\
\bottomrule
\end{tabular}
\caption{Pass rate (\%) by condition and task complexity. Paired exact McNemar:
gemma EN\,vs\,BN $p=8.2{\times}10^{-5}$ ($b{=}41$, $c{=}12$);
gpt-oss EN\,vs\,BN $p=0.043$ ($b{=}23$, $c{=}40$).}
\label{tab:main}
\end{table}

The aggregate conclusion is therefore not ``Bangla is harder''. The language manipulation is
identical; only the model differs. Any claim of the form ``model performance degrades in
language $L$'' is under-specified unless the model is named.

\paragraph{Rankings do not transfer across languages.}
This has a direct practical consequence (Table~\ref{tab:rank}). In English, \texttt{gemma4:31b}
beats \texttt{gpt-oss:20b} decisively: 67 subtasks are solved only by gemma against 15 only by
gpt-oss ($p=5.3\times10^{-9}$). In Bangla the same comparison yields 33 against 27---
statistically indistinguishable ($p=0.52$). A practitioner choosing a model from an English
leaderboard would believe the first model is far better; in Bangla that advantage is not
measurable.

\begin{table}[t]
\centering\small
\begin{tabular}{lccc}
\toprule
Comparison & $b$ & $c$ & exact $p$ \\
\midrule
gemma vs.\ gpt-oss (EN) & 67 & 15 & $5.3{\times}10^{-9}$ \\
gemma vs.\ gpt-oss (BN) & 33 & 27 & $0.52$ \\
\bottomrule
\end{tabular}
\caption{The model gap is overwhelming in English and vanishes in Bangla. $b$: solved only by
gemma; $c$: solved only by gpt-oss.}
\label{tab:rank}
\end{table}

\subsection{RQ2: Complexity gates the language effect}

Table~\ref{tab:main} also shows that the language effect is not uniform across difficulty. For
\texttt{gemma4:31b} the Bangla penalty is largest on one-app tasks ($-21.5$ points) and
\emph{exactly zero} on three-app tasks ($16.1\%$ in both conditions).

The explanation is a floor effect. Three-app subtasks require long multi-application
workflows---read an image, OCR it, summarise, split the result, send mail---and both models fail
the large majority of them regardless of language. Where a model has no competence to lose,
language cannot take any away. The corollary matters for benchmark design: a cross-lingual
evaluation run only on hard tasks will under-report the language effect, because the hardest
items are least sensitive to it.

\subsection{RQ3: Failure mode determines the sign}

Why do the two models move in opposite directions under an identical manipulation? Because they
fail for different reasons (Table~\ref{tab:modes}).

\begin{table}[t]
\centering\small
\begin{tabular}{lcccc}
\toprule
 & \multicolumn{2}{c}{gemma4:31b} & \multicolumn{2}{c}{gpt-oss:20b} \\
\cmidrule(lr){2-3}\cmidrule(lr){4-5}
Signal & EN & BN & EN & BN \\
\midrule
Failures            & 168 & 197 & 220 & 203 \\
\;\;malformed action & 53 & 23 & \textbf{212} & \textbf{186} \\
\;\;ended stuck      & 3 & 5 & 30 & 14 \\
\;\;out of steps     & 46 & 39 & 81 & 58 \\
\;\;finished, failed & 118 & 152 & 108 & 130 \\
\midrule
Mean steps/episode   & 10.9 & 10.5 & 13.3 & 12.1 \\
Bangla-script answer & 0 & \textbf{51} & 0 & 7 \\
\bottomrule
\end{tabular}
\caption{Failure composition over all 300 subtasks per condition. The three terminal states
(stuck / out of steps / finished-but-failed) are mutually exclusive and sum to the failure count
together with one unsaved episode per condition; \emph{malformed action} is a cross-cutting flag
counting failures in which at least one action was rejected, so it overlaps the others. gpt-oss
is dominated by action-format failures; gemma by episodes that finish and still fail.}
\label{tab:modes}
\end{table}

\texttt{gpt-oss:20b} is \textbf{format-bound}. $212$ of its $220$ English failures ($96\%$)
contain at least one action rejected as malformed. Its reasoning lives in a separate channel,
and in English it frequently spills prose into the action slot, producing a rejected action,
then another, until the episode is consumed. Under Bangla instructions this leakage falls:
malformed-action failures drop to $186$, stuck episodes halve ($30\rightarrow14$), budget
exhaustion falls ($81\rightarrow58$), and episodes shorten ($13.3\rightarrow12.1$ steps). The
model is not understanding the task better; it is \emph{talking less} into the action channel,
and it therefore completes more episodes. Its gain in Bangla is a format effect, not a
comprehension effect.

\texttt{gemma4:31b} is \textbf{content-bound}. It emits valid actions (only $53$ of $168$
English failures involve a malformed action) and its dominant failure is an episode that runs to
completion and is then rejected on content---$118$ in English rising to $152$ in Bangla. This is
the mode that a language change should affect, and it does.

This yields a simple predictive rule, and the central claim of this paper:

\begin{quote}\itshape
The sign of an agent's cross-lingual effect follows its dominant failure mode. Content-bound
agents lose under a language shift; format-bound agents may gain, because the shift suppresses
the verbosity that corrupts their actions.
\end{quote}

\subsection{RQ4: Much of the measured penalty belongs to the evaluator}

The last row of Table~\ref{tab:modes} isolates a failure mode that \emph{cannot} occur in
English. In $51$ of its $197$ Bangla failures ($26\%$), \texttt{gemma4:31b} terminated by
calling \texttt{finish\_task} with an answer written in Bangla script. For $49$ of those
subtasks the failing check is \texttt{evaluate\_contain}, which tests whether an English keyword
occurs in the output.

In other words: the agent was asked a question in Bangla, answered in Bangla, and was marked
wrong because the checker was written in English. This is not an agent capability failure. It is
a property of the measuring instrument.

Table~\ref{tab:bounds} reports the sensitivity. Under the strict (as-shipped) checkers
\texttt{gemma4:31b} scores $34.3\%$ in Bangla. If every Bangla-script answer were credited---a
deliberate upper bound, since we do not verify each one semantically---it would score $51.3\%$,
which is \emph{above} its own English score of $44.0\%$. The true value lies between the two
bounds, and we cannot determine it with English-only checkers. The honest conclusion is
therefore stronger and more useful than the naive one: \textbf{we cannot claim from these
numbers that Bangla degrades this model's task competence, because the instrument that produced
the number is itself language-biased.}

\begin{table}[t]
\centering\small
\begin{tabular}{lccc}
\toprule
Condition & Strict & Lenient bound & EN reference \\
\midrule
gemma4:31b BN  & 34.3 & $\le 51.3$ & 44.0 \\
gpt-oss:20b BN & 32.3 & $\le 34.7$ & 26.7 \\
\bottomrule
\end{tabular}
\caption{Sensitivity of the Bangla score to evaluator language (\%). The lenient bound credits
every Bangla-script answer. For gemma the bound brackets---and exceeds---the English score.}
\label{tab:bounds}
\end{table}

Notably the confound is itself model-dependent: \texttt{gpt-oss:20b} answered in Bangla only $7$
times, because it remains English-anchored even under Bangla instructions. A checker bias thus
penalises precisely the model that adapts to the user's language---exactly the behaviour one
would want to reward.

\section{Discussion}

\paragraph{``Multilingual degradation'' is not one phenomenon.} Our two models received an
identical manipulation and moved in opposite directions, with the sign predicted by their
failure composition. Reporting a single aggregate delta per language hides this. We recommend
that cross-lingual agent evaluations report failure-mode composition alongside the score; it is
both more informative and, as we show, necessary to interpret the sign.

\paragraph{Leaderboards are language-scoped artefacts.} The ranking between our two models is
decisive in English and absent in Bangla. Deployment decisions for non-English users should not
be read off English leaderboards.

\paragraph{Evaluation harnesses need language-aware verification.} Execution-based checks are
usually presented as more robust than surface matching. They are---for file state and cell
values. But a large fraction of office tasks are questions whose answer is checked by keyword,
and those checks silently encode a language. Multilingual agent benchmarks should either verify
answers semantically, accept the task language, or separate ``artefact-state'' checks from
``answer-text'' checks and report them apart. Our Table~\ref{tab:bounds} is the minimum
diagnostic we would want to see reported.

\paragraph{Interface language versus instruction language.} Keeping the tool schema in English
was a methodological necessity, but it reflects a real deployment regime and it isolates a
useful finding: an agent can be instructed in one language while acting through an
English-symbol API, and the main cost is not symbol confusion but answer-language mismatch at
the boundary where the agent reports back to the user.

\section{Conclusion}

We built a strictly paired English/Bangla agent evaluation over \textsc{OfficeBench} and ran
1{,}200 episodes across two open-weight models. Task language changes agent success, but the
direction depends on the model's dominant failure mode: a content-bound model loses $9.7$
points, while a format-bound model gains $5.6$. Model rankings that are decisive in English can
vanish in another language. And a quarter of one model's Bangla failures are artefacts of
English-only checkers, enough to bracket its true Bangla score above its English score. Language
matters for LLM agents---but how it matters is a property of the model and of the evaluator, not
of the language alone.

\section*{Limitations}

\paragraph{One non-English language.} We study Bangla only. Bangla is a useful probe---widely
spoken yet under-represented---but our directional claim (sign follows failure mode) should be
tested on typologically different and higher-resource languages before being treated as general.

\paragraph{Machine-translated instructions.} Task instructions were machine-translated and
cached; the system prompt was authored natively. Translation noise is therefore part of the
Bangla condition, and some failures we code as content errors may be translation artefacts. We
mitigated this by using one fixed translation per subtask for both models, so the comparison
\emph{between} models is unaffected, but the absolute Bangla scores should be read as a lower
bound on achievable performance.

\paragraph{Two models.} Two models cannot establish a general law. They were chosen to differ in
kind (conventional versus separate-reasoning-channel), which is what exposes the mechanism, but
the mechanism now needs replication across more systems. An evaluation of a third model in the
same size class was in progress at submission time and is not reported here.

\paragraph{Decoding is not perfectly reproducible.} We decode at temperature $0$, but served
inference is not bit-deterministic: batching and floating-point non-associativity can change an
argmax when two tokens are near-tied, and a single divergence early in an episode can flip the
outcome. We did not re-run conditions to estimate this variance directly, so part of the smaller
effect we report (the $5.6$-point gpt-oss gain, $p=0.043$) may be run-to-run noise. The larger
gemma effect ($p<0.001$) and the model-gap collapse are well outside plausible noise, but a
replication study would strengthen the weaker claim. We regard this as the most important
missing control.

\paragraph{The lenient bound is an upper bound.} We did not semantically verify each of the 51
Bangla-script answers; the bound credits them all. The true Bangla score lies strictly between
the two bounds and we do not claim to have located it.

\paragraph{Infrastructure effects.} Episodes ran on a memory-constrained host. A small number of
episodes (at most one per condition in the reported runs) failed to persist their output and are
counted as failures, and a per-episode wall-clock cap occasionally terminated a slow episode.
These affect all conditions approximately equally and are too few to alter any reported
comparison, but they are not zero.

\section*{Ethics Statement}

This work evaluates publicly available open-weight models on a public benchmark in a sandboxed
container; no human subjects and no private data are involved. Our main finding has a direct
equity dimension: evaluation harnesses written in English can penalise an agent for answering in
the user's language, which systematically understates agent quality for non-English users and
rewards models that ignore the user's language. We release our bilingual protocol, trajectories
and failure codes so that this bias can be measured rather than inherited. The translated task
set was produced by machine translation and may contain errors; it is intended for evaluation
research, not as a curated Bangla language resource.

\bibliography{refs}
% acl.sty already sets \bibliographystyle{acl_natbib}; setting it again makes BibTeX error.

\appendix

\section{Reproducibility details}
\label{app:repro}

Each episode runs in a fresh container reset from a pinned image, with a 20-step budget and a
single seeded file system per subtask. Decoding is fixed at $\mathrm{temperature}=0$,
$\mathrm{top\_p}=1$. The harness records, per step, the full prompt, the model reply, the parsed
action and the environment observation; all analyses in this paper are computed from those logs
and from the benchmark's own scorer output, not from intermediate summaries.

\section{Failure code book}
\label{app:codes}

\begin{table}[h]
\centering\small
\begin{tabular}{lp{5.1cm}}
\toprule
Code & Meaning \\
\midrule
BM  & Benchmark defect: a check or reference asks for something the task does not state \\
WC  & Wrong content: output exists, a value is wrong \\
SB  & Out of steps: hit the 20-step budget \\
PF  & Premature finish: declared done with a required output missing \\
TR  & Task read differently: a different operation from the one asked \\
FMT & Action format: could not produce a valid action \\
ER  & Empty result taken as ``nothing exists'' \\
\midrule
\multicolumn{2}{l}{\emph{Bangla-condition only}} \\
AL  & Answer language: correct answer in Bangla script, English check fails \\
NC  & Name case: a name rewritten from Bangla script with wrong case \\
WD  & Word choice when translating a literal back to English \\
TX  & Translation changed the task or a literal \\
\bottomrule
\end{tabular}
\caption{Failure codes. AL--TX cannot occur in the English condition and are the
language-specific modes discussed in Section~4.4.}
\label{tab:codebook}
\end{table}

\end{document}
